\documentclass[a4paper, serif, 10pt]{moderncv}

%% moderncv
% color and style
\moderncvcolor{black}
\moderncvstyle{banking}

% renew moderncv command to adapt for CJK colon
\renewcommand*{\cvitem}[3][.25em]{%
  \ifstrempty{#2}{}{\hintstyle{#2}：}{#3}%
  \par\addvspace{#1}}

\renewcommand*{\cvitemwithcomment}[4][.25em]{%
  \savebox{\cvitemwithcommentmainbox}{\ifstrempty{#2}{}{\hintstyle{#2}：}#3}%
  \setlength{\cvitemwithcommentmainlength}{\widthof{\usebox{\cvitemwithcommentmainbox}}}%
  \setlength{\cvitemwithcommentcommentlength}{\maincolumnwidth-\separatorcolumnwidth-\cvitemwithcommentmainlength}%
  \begin{minipage}[t]{\cvitemwithcommentmainlength}\usebox{\cvitemwithcommentmainbox}\end{minipage}%
  \hfill% fill of \separatorcolumnwidth
  \begin{minipage}[t]{\cvitemwithcommentcommentlength}\raggedleft\small\itshape#4\end{minipage}%
  \par\addvspace{#1}}

%% page layout/margins
\usepackage[top=2.5cm, bottom=2.5cm, left=1.5cm, right=1.5cm]{geometry}

\nopagenumbers{}

%% Babel config for English language
\usepackage[english]{babel}

%% fontspec
\usepackage{fontspec}

\IfFontExistsTF{Linux Libertine}{
  \setmainfont[Ligatures={TeX, Common}, Numbers=Lining, ItalicFont=Linux Libertine]{Linux Libertine}
}{}
\IfFontExistsTF{Linux Libertine O}{
  \setmainfont[Ligatures={TeX, Common}, Numbers=Lining, ItalicFont=Linux Libertine O]{Linux Libertine O}
}{}

%% CTeX
% CJK support, used to show CJK characters in the resume
%
% - heading=false: disable ctex heading
% - punct=kaiming: use kaiming punctuations styles for CJK
% - scheme=plain: use plain scheme, do not override \normalsize font size
% - space=auto: space settings for CJK characters
%
% CJK fonts are resolved per locale language, e.g. Noto Serif/Sans CJK JP
% for Japanese resumes, with fallback to the next region if not installed
%
% ref:
% - http://ctan.mirrorcatalogs.com/language/chinese/ctex/ctex.pdf
\usepackage[UTF8, heading=false, punct=kaiming, scheme=plain, space=auto]{ctex}

\IfFontExistsTF{Noto Serif CJK JP}{
  \setCJKmainfont{Noto Serif CJK JP}
}{}
\IfFontExistsTF{Noto Sans CJK JP}{
  \setCJKsansfont{Noto Sans CJK JP}
}{}

%% line spacing
\usepackage{setspace}
\setstretch{1.125}

%% URL styling - use normal text instead of monospace
\urlstyle{same}

%% Patch \httplink and \httpslink to handle full URLs
% The original moderncv macros always prepend http:// or https:// to URLs.
% This causes issues when the URL already has a protocol (e.g., https://example.com)
% resulting in malformed URLs like https://https://example.com.
% This patch checks if the URL already contains "://" and uses it directly if so.
% Note: We use \str_set:Nx (x-expansion) to expand macros like \@homepage first.
\ExplSyntaxOn
\str_new:N \g_yamlresume_url_str

\RenewDocumentCommand{\httpslink}{O{}m}{
  \str_set:Nx \g_yamlresume_url_str {#2}
  \str_if_in:NnTF \g_yamlresume_url_str {://}
    { \url{#2} }
    { \url{https://#2} }
}

\RenewDocumentCommand{\httplink}{O{}m}{
  \str_set:Nx \g_yamlresume_url_str {#2}
  \str_if_in:NnTF \g_yamlresume_url_str {://}
    { \url{#2} }
    { \url{http://#2} }
}
\ExplSyntaxOff

\name{佐藤 美咲}{}
\title{法人SaaS営業のアカウントエグゼクティブ}
\phone[mobile]{+81 90-1234-5678}
\email{misaki.sato@example.co.jp}
\homepage{https://misaki-sato.example.com}

\address{日本、東京都、千代田区、東京都千代田区丸の内1-1-1、100-0005}{}{}

\extrainfo{{\small \faLinkedin}\ \href{https://www.linkedin.com/in/misaki-sato}{@misaki-sato} {} {} {} • {} {} {} 
{\small \faTwitter}\ \href{https://twitter.com/misaki\_sato\_sales}{@misaki\_sato\_sales}}

\begin{document}

\maketitle

\section{基本情報}

\cvline{}{\begin{itemize}
\item 法人向けSaaS営業で7年以上の経験を持つアカウントエグゼクティブ
\item エンタープライズ企業を中心に、新規開拓からクロージング、顧客成功まで一貫して担当
\item 年間売上目標を120\%達成し、複数回の最優秀営業賞を受賞
\item データに基づく商談設計と、経営層との関係構築を得意とする
\end{itemize}}
\section{学歴}

\cventry{2014年4月～2018年3月}
        {学士、経済学部経済学科、成績：3.7}
        {慶應義塾大学}
        {\url{https://www.keio.ac.jp/}}
        {}
        {\begin{itemize}
\item 経済学と経営学を中心に、市場分析と企業戦略の基礎を修得
\item ゼミではB2Bマーケティングを研究し、企業への提案を実施
\item 交換留学で英語によるディスカッションとプレゼンテーションを経験
\end{itemize}
\textbf{コース}：マーケティング論、消費者行動論、企業戦略論、会計学、統計学、交渉学、ビジネス英語、組織行動論}
\section{職歴}

\cventry{2021年4月～現在}
        {シニアアカウントエグゼクティブ}
        {株式会社クラウドセールス}
        {\url{https://www.cloudsales.example.co.jp}}
        {}
        {\begin{itemize}
\item エンタープライズ向けSaaSの新規開拓営業を担当し、年間売上目標を120\%達成
\item 経営層への提案からPoC、契約交渉までをリードし、大型案件を複数成約
\item 顧客成功チームと連携し、オンボーディングとアップセルを推進
\item 若手営業のメンターとして商談スキル向上を支援
\end{itemize}
\textbf{キーワード}：SaaS、法人営業、エンタープライズ、新規開拓、クロージング、顧客成功}

\cventry{2018年4月～2021年3月}
        {アカウントエグゼクティブ}
        {株式会社データブリッジ}
        {\url{https://www.databridge.example.co.jp}}
        {}
        {\begin{itemize}
\item 中堅・中小企業向けにデータ分析ツールの営業を担当
\item インサイドセールスとフィールドセールスを兼任し、商談数を3倍に拡大
\item CRMを活用したパイプライン管理と営業プロセスの改善を推進
\item 顧客の課題ヒアリングに基づく提案で、解約率を15\%削減
\end{itemize}
\textbf{キーワード}：データ分析、CRM、パイプライン管理、提案営業、課題解決}
\section{言語}

\cvline{日本語}{ネイティブ・母国語レベル \hfill \textbf{キーワード}：ビジネス日本語、母国語}
\cvline{英語}{完全な実務レベル \hfill \textbf{キーワード}：TOEIC 950、ビジネス交渉}
\cvline{中国語}{限定的な実務レベル \hfill \textbf{キーワード}：HSK 4級、日常会話}
\section{スキル}

\cvline{法人営業}{エキスパート \hfill \textbf{キーワード}：新規開拓、商談設計、クロージング、顧客関係構築、アップセル}
\cvline{SaaSソリューション}{上級 \hfill \textbf{キーワード}：Salesforce、HubSpot、KPI管理、データ分析、提案資料作成}
\cvline{プレゼンテーション}{上級 \hfill \textbf{キーワード}：経営層向け提案、デモ、交渉、ファシリテーション}
\cvline{プロジェクトマネジメント}{中級 \hfill \textbf{キーワード}：ステークホルダー管理、スケジュール管理、リスク管理}
\section{受賞・表彰}

\cventry{2023年12月}
        {株式会社クラウドセールス}
        {年間最優秀営業賞}
        {}
        {}
        {年間売上目標を120\%達成し、新規顧客獲得数が全社1位。エンタープライズ案件の成約額でも最高記録を更新。}

\cventry{2022年6月}
        {株式会社クラウドセールス}
        {四半期MVP}
        {}
        {}
        {第2四半期において、パイプライン創出額と成約率が全営業部門でトップ。}
\section{資格・免状}

\cventry{2022年6月}
        {Salesforce}
        {Salesforce認定アドミニストレーター}
        {\url{https://trailhead.salesforce.com/credentials/administrator}}
        {}
        {}

\cventry{2020年11月}
        {IIBC}
        {TOEIC Listening \& Reading 950}
        {\url{https://www.iibc-global.org/english/toeic.html}}
        {}
        {}
\section{出版・著作}

\cventry{2022年8月}
        {法人SaaS営業における顧客成功指標の設計}
        {日本セールス研究ジャーナル}
        {\url{https://example.com/journal/sales-cs}}
        {}
        {\begin{itemize}
\item 顧客成功指標と営業KPIの連携方法を提案
\item 解約率低下とアップセル率向上の実践事例を分析
\end{itemize}}
\section{推薦人}

\cventry{\emaillink[ichiro.tanaka@cloudsales.example.co.jp]{ichiro.tanaka@cloudsales.example.co.jp}}
        {株式会社クラウドセールス 営業部長}
        {田中 一郎}
        {+81 3-1234-5678}
        {}
        {佐藤美咲は卓越した営業スキルと顧客志向を持ち、組織に大きな貢献をもたらしました。どんなチームでも即戦力として活躍できる人材です。}
\section{プロジェクト}

\cventry{2023年1月～2023年8月}
        {大手製造業向けにCRMを導入し、営業プロセスを標準化するプロジェクト}
        {エンタープライズ向けCRM導入プロジェクト}
        {\url{https://www.cloudsales.example.co.jp/projects/crm}}
        {}
        {\begin{itemize}
\item 顧客の営業部門と連携し、要件定義から運用定着までをリード
\item 導入後に商談サイクルを30\%短縮し、受注率を15\%向上
\item チェンジマネジメントを推進し、現場の利用率を90\%以上に向上
\end{itemize}
\textbf{キーワード}：CRM、要件定義、チェンジマネジメント、営業プロセス}
\section{興味・関心}

\cvline{スポーツ}{テニス、ランニング、登山}
\cvline{読書}{ビジネス書、歴史小説、経済学}
\section{ボランティア活動}

\cventry{2022年4月～現在}
        {メンター}
        {一般社団法人次世代リーダー育成協会}
        {\url{https://example.org/next-leaders}}
        {}
        {\begin{itemize}
\item 大学生向けに営業スキルとキャリア形成のワークショップを開催
\item 年間10回のメンタリングセッションを担当
\end{itemize}}
    
\end{document}